\documentclass[11pt]{article}

\usepackage{amsmath,amssymb,amsthm,mathtools}
\usepackage[margin=1in]{geometry}
\usepackage{longtable}
\usepackage{url}

\newtheorem{theorem}{Theorem}[section]
\newtheorem{lemma}[theorem]{Lemma}
\newtheorem{corollary}[theorem]{Corollary}

\newcommand{\dd}{\,d}

\begin{document}

\section{Purpose and arbitrary-fibre setup}

For a finite graph $H$, put
\[
 \Gamma_H(W)=t(H,W)-\Phi_H(p(W)),
 \qquad
 \Phi_H(p)=(1-p)^{v(H)}
 \chi_H\!\left(\frac1{1-p}\right).
\]
The constant two-block hub theorem leaves open nonconstant fibres from the
exceptional set into the complete core.  This note first removes that
restriction with an adaptive exact radius and then proves a uniform radius
with no restriction on either the cut or exceptional-square kernel.

Let $(\Omega,\mu)$ be an arbitrary probability space and split it into
measurable sets $A,B$ of masses $s,1-s$, where $0<s<1$.  Normalize the
restrictions of $\mu$ to probability measures $\mu_A,\mu_B$.  Let $W$ be any
graphon satisfying only
\begin{equation}
 W(x,y)=1\quad\text{for almost every }(x,y)\in B^2.
 \label{eq:complete-core}
\end{equation}
Both $W|_{A\times B}$ and $W|_{A^2}$ are otherwise arbitrary measurable
kernels.  Put $V=1-W$ and, for almost every $z\in A$, define the normalized
cross-complement degree
\begin{equation}
 a(z)=\int_B V(z,y)\dd\mu_B(y),
 \qquad
 E_2=\int_Aa(z)^2\dd\mu_A(z).
 \label{eq:cross-energy}
\end{equation}
The theorem proves positivity whenever $s\leq\rho_HE_2$.  In particular, if
$E_2$ stays bounded below while $s\downarrow0$, it applies uniformly even
though vertices in $A$ retain complement degree bounded away from zero.


\section{The exact first layer and a scalar coercivity lemma}

Write $n=v(H)$, $m=e(H)$, and $d_v=\deg_H(v)$.  Since all integrands are
nonnegative, retain only maps sending zero or one vertices of $H$ into $A$.
The zero-vertex contribution is $(1-s)^n$.  If exactly $v$ maps to a point
$z\in A$, every neighbor of $v$ is an independent $\mu_B$-variable and all
edges among the remaining vertices have value one by
\eqref{eq:complete-core}.  Thus
\begin{equation}
 t(H,W)\geq(1-s)^n+s(1-s)^{n-1}
 \sum_{v\in V(H)}\int_A(1-a(z))^{d_v}\dd\mu_A(z).
 \label{eq:first-layer}
\end{equation}
This identity retains the full arbitrary fibre through its exact degree
$a(z)$; no product-of-means replacement is made.

\begin{lemma}[Degree-sequence coercivity]
\label{lem:scalar}
Suppose $H$ has no isolated vertices.  For every $u\in[0,1]$,
\begin{equation}
 2mu-n+\sum_{v\in V(H)}(1-u)^{d_v}
 \geq(2m-n)u^2.
 \label{eq:scalar}
\end{equation}
\end{lemma}

\begin{proof}
For every integer $d\geq1$,
\begin{equation}
 du-1+(1-u)^d\geq(d-1)u^2.
 \label{eq:one-degree}
\end{equation}
To see this, put $t=1-u$.  The difference between the two sides is
\[
 t\bigl(t^{d-1}+(d-2)-(d-1)t\bigr)\geq0,
\]
because the convex function $t^{d-1}$ lies above its tangent at $t=1$.
Sum \eqref{eq:one-degree} over the vertices and use
$\sum_vd_v=2m$ and $\sum_v(d_v-1)=2m-n$.
\end{proof}


\section{Uniform target remainder and the graphon theorem}

Write the polynomial continuation at high density as
\begin{equation}
 \Psi_H(q):=\Phi_H(1-q)
 =1-mq+\sum_{j=2}^nb_jq^j,
 \qquad
 C_H=\sum_{j=2}^n|b_j|.
 \label{eq:target-remainder}
\end{equation}
The coefficient of $q$ is $-m$ because the two leading coefficients of the
chromatic polynomial are $1,-m$.

\begin{theorem}[Arbitrary-fibre persistent-hub stability]
\label{thm:hub}
Let $H$ be a finite graph with no isolated vertices and $2m>n$.  Define
\begin{equation}
 K_H=n(n-1)+2m+4C_H,
 \qquad
 \rho_H=\frac{2m-n}{K_H}.
 \label{eq:constants}
\end{equation}
For every graphon satisfying \eqref{eq:complete-core},
\begin{equation}
 \Gamma_H(W)\geq
 s(2m-n)E_2-K_Hs^2.
 \label{eq:main-bound}
\end{equation}
Consequently,
\begin{equation}
 0\leq s\leq\rho_HE_2
 \quad\Longrightarrow\quad
 t(H,W)\geq\Phi_H(p(W)).
 \label{eq:adaptive-radius}
\end{equation}
The inequality is strict when $0<s<\rho_HE_2$.
\end{theorem}

\begin{proof}
Let
\[
 \bar a=\int_Aa\dd\mu_A,
 \qquad
 \bar r=\int_{A^2}V\dd\mu_A^2.
\]
Direct integration gives
\begin{equation}
 q=1-p(W)=2s(1-s)\bar a+s^2\bar r.
 \label{eq:q}
\end{equation}
In particular
\[
 2s(1-s)\bar a\leq q\leq2s-s^2\leq2s.
\]
Since $(1-s)^k\geq1-ks$ for an integer $k\geq1$, and the sum in
\eqref{eq:first-layer} is at most $n$, that formula implies
\begin{align}
 t(H,W)
 &\geq1+s\left(
 -n+\sum_v\int_A(1-a)^{d_v}\dd\mu_A
 \right)-n(n-1)s^2.
 \label{eq:hom-lower}
\end{align}
On the other hand, \eqref{eq:target-remainder} and \eqref{eq:q} give
\begin{align}
 \Psi_H(q)
 &\leq1-mq+C_Hq^2\\
 &\leq1-2ms\bar a+(2m+4C_H)s^2.
 \label{eq:target-upper}
\end{align}
Subtract \eqref{eq:target-upper} from \eqref{eq:hom-lower}.  Applying
Lemma~\ref{lem:scalar} pointwise with $u=a(z)$ gives
\[
 \Gamma_H(W)
 \geq s(2m-n)\int_Aa^2\dd\mu_A
 -\bigl(n(n-1)+2m+4C_H\bigr)s^2,
\]
which is \eqref{eq:main-bound}.  The two conclusions follow immediately.
All estimates remain valid when the internal kernel $V|_{A^2}$ is arbitrary;
it enters only through the favorable lower bound on $q$ and the universal
upper bound $q\leq2s$.
\end{proof}


\section{A sharp union-bound improvement and the complete cut-only regime}

The preceding theorem treats an arbitrary exceptional-square kernel but loses
an $O(s^2)$ term even when that kernel and the cross defect both vanish.  A
more faithful comparison removes that loss and closes the entire cut-only
transition, including $E_2=o(s)$.

\begin{theorem}[Exact one-exceptional-vertex gain]
\label{thm:union-improvement}
In the setup of Theorem~\ref{thm:hub}, let $q=1-p(W)$.  Then
\begin{equation}
 \Gamma_H(W)\geq
 s(1-s)^{n-1}(2m-n)E_2-C_Hq^2.
 \label{eq:sharp-bound}
\end{equation}
This bound still permits an arbitrary kernel on $A^2$.
\end{theorem}

\begin{proof}
For numbers $v_e\in[0,1]$,
\[
 \prod_e(1-v_e)\geq1-\sum_ev_e.
\]
Apply this pointwise with the complement factors on the edges of $H$ and
integrate.  The integral of the right side is exactly $1-mq$.  The omitted
nonnegative remainder may be retained on any selected assignments.

Consider assignments for which exactly one graph vertex $v$ lands in $A$.
Conditional on its value $z\in A$, every other vertex lies in $B$.  All
nonincident edges have graphon value one, while the $d_v$ neighbor variables
are independent under the product measure.  Hence the exact product and its
linear union bound integrate respectively to
\[
 (1-a(z))^{d_v},
 \qquad1-d_va(z).
\]
The total gain on all such assignments is therefore
\[
 s(1-s)^{n-1}\sum_v\int_A
 \bigl((1-a)^{d_v}-1+d_va\bigr)\dd\mu_A.
\]
The proof of Lemma~\ref{lem:scalar}, before summing its linear terms, shows
that the summand is at least $(d_v-1)a^2$.  Summing gives
\begin{equation}
 t(H,W)\geq1-mq+s(1-s)^{n-1}(2m-n)E_2.
 \label{eq:improved-union}
\end{equation}
Finally, \eqref{eq:target-remainder} and $0\leq q\leq1$ imply
\[
 \Psi_H(q)\leq1-mq+C_Hq^2.
\]
Subtract this from \eqref{eq:improved-union}.
\end{proof}

\begin{corollary}[Arbitrary crossing-cut fibres]
\label{cor:cut-only}
Assume in addition that $W=1$ almost everywhere on $A^2$; equivalently, the
complement is supported only on $A\times B$ and $B\times A$.  Put
\begin{equation}
 \sigma_H=
 \frac{2m-n}{(2m-n)(n-3)+4C_H}.
 \label{eq:fixed-radius}
\end{equation}
Then
\[
 0\leq s\leq\sigma_H
 \quad\Longrightarrow\quad
 t(H,W)\geq\Phi_H(p(W))
\]
for every measurable cut kernel.  No lower bound on $E_2$ is required.
Equality occurs when the cut defect vanishes; for $0<s<\sigma_H$ the
inequality is strict whenever $E_2>0$.
\end{corollary}

\begin{proof}
Here
\[
 q=2s(1-s)\bar a,
 \qquad \bar a=\int_Aa\dd\mu_A,
\]
and Cauchy--Schwarz gives $q^2\leq4s^2(1-s)^2E_2$.
Theorem~\ref{thm:union-improvement} therefore yields
\[
 \Gamma_H(W)\geq
 s(1-s)^2E_2
 \left((2m-n)(1-s)^{n-3}-4C_Hs\right).
\]
Bernoulli's inequality gives
$(1-s)^{n-3}\geq1-(n-3)s$.  The bracket is consequently at least
\[
 (2m-n)-\bigl((2m-n)(n-3)+4C_H\bigr)s,
\]
which is nonnegative precisely throughout the claimed rational interval.
The equality statements follow from the same display.
\end{proof}


\section{The fully arbitrary single-hub theorem}

The cut-only corollary and the internal maximum-degree theorem combine without
any relation between the two kernels.  For
\[
 \Psi_H(q)=1-mq+R_H(q),
 \qquad R_H(q)=\sum_{j=2}^nb_jq^j,
\]
put
\begin{equation}
 D_H=\sum_{j=2}^n j|b_j|.
 \label{eq:derivative-norm}
\end{equation}
Let
\[
 g_H=\sum_{v\in V(H)}\binom{d_v}{2}-N_3(H).
\]
Use the graph constants $\Lambda_H,\Omega_H$ from the maximum-degree
complement theorem: for any complement of maximum degree $\Delta$,
\begin{equation}
 \Gamma_H(1-V)\geq
 (g_H-\Lambda_H\Delta)X_V
 +(N_3(H)-\Omega_H\Delta)Y_V.
 \label{eq:maximum-degree-input}
\end{equation}

\begin{lemma}[Superadditivity of the union-bound remainder]
\label{lem:superadditive}
For two disjoint collections of numbers in $[0,1]$, put
\[
 F(\mathcal E)=\prod_{e\in\mathcal E}(1-v_e)-1+
 \sum_{e\in\mathcal E}v_e.
\]
Then
\begin{equation}
 F(\mathcal E_0\cup\mathcal E_1)
 =F(\mathcal E_0)+F(\mathcal E_1)
 +(1-P_0)(1-P_1)\geq F(\mathcal E_0)+F(\mathcal E_1),
 \label{eq:superadditive}
\end{equation}
where $P_i=\prod_{e\in\mathcal E_i}(1-v_e)$.
\end{lemma}

\begin{proof}
Expand the right side of the identity.  Its value is
$P_0P_1-1+\sum_{\mathcal E_0\cup\mathcal E_1}v_e$.  The final term is
nonnegative because $0\leq P_i\leq1$.
\end{proof}

\begin{theorem}[Complete-core, fully arbitrary single hub]
\label{thm:complete-core}
Suppose $H$ has no isolated vertices and $W=1$ almost everywhere on $B^2$, as in
\eqref{eq:complete-core}, with no restriction on either $W|_{A\times B}$ or
$W|_{A^2}$.  Put $a_H=2m-n$ and
\begin{equation}
 \tau_H=\min\left\{
 \frac12,
 \frac{g_H}{2\Lambda_H},
 \frac{N_3(H)}{\Omega_H},
 \frac{a_Hg_H}
 {a_Hg_H(n-1)+4D_Hg_H+8D_H^2}
 \right\}.
 \label{eq:complete-core-radius}
\end{equation}
If $a_H,g_H,N_3(H)>0$ and $0\leq s\leq\tau_H$, then
\[
 t(H,W)\geq\Phi_H(p(W)).
\]
For $0<s<\tau_H$, the inequality is strict unless $W=1$ almost everywhere.
\end{theorem}

\begin{proof}
Split the complement into its internal and cut parts,
\[
 V_0=V\mathbf1_{A^2},
 \qquad V_1=V-V_0,
 \qquad q_i=\int V_i.
\]
Write $W_i=1-V_i$.  Applying Lemma~\ref{lem:superadditive} pointwise to the
edge defects and integrating gives
\begin{equation}
 \Gamma_H(W)\geq
 \Gamma_H(W_0)+s(1-s)^{n-1}a_HE_2
 -\bigl(R_H(q_0+q_1)-R_H(q_0)\bigr).
 \label{eq:combined-gap}
\end{equation}
Indeed, the integrated internal remainder is
$t(H,W_0)-1+mq_0$, while the integrated cut remainder is at least the exact
one-exceptional-vertex gain used in
Theorem~\ref{thm:union-improvement}; the linear target terms cancel.

Let
\[
 \bar r=\int_{A^2}V\dd\mu_A^2,
 \qquad \bar a=\int_Aa\dd\mu_A.
\]
Then $q_0=s^2\bar r$ and
$q_1=2s(1-s)\bar a\leq2s\sqrt{E_2}$.  Since
\[
 0\leq q^j-q_0^j\leq jq_1q^{j-1}\leq jq_1(q_0+q_1)
 \quad(j\geq2),
\]
the target difference in \eqref{eq:combined-gap} has absolute value at most
\begin{equation}
 D_H(q_0q_1+q_1^2)
 \leq D_H\bigl(2s^3\bar r\sqrt{E_2}+4s^2E_2\bigr).
 \label{eq:target-difference}
\end{equation}

The complement $V_0$ has maximum degree at most $s$.  Under the second and
third restrictions in \eqref{eq:complete-core-radius}, the input
\eqref{eq:maximum-degree-input} gives
\[
 \Gamma_H(W_0)\geq\frac{g_H}{2}X_{V_0}.
\]
If $R=V|_{A^2}$ is viewed on the normalized probability space $A$, then
\[
 X_{V_0}=s^3t(P_3,R)-s^4\bar r^2
 \geq s^3(1-s)\bar r^2,
\]
where $t(P_3,R)=\int d_R^2\geq\bar r^2$ by Cauchy--Schwarz.  Consequently
\begin{align}
 \Gamma_H(W)
 &\geq \frac{g_H}{2}s^3(1-s)\bar r^2
 +a_Hs(1-s)^{n-1}E_2 \\
 &\quad-D_H\bigl(2s^3\bar r\sqrt{E_2}+4s^2E_2\bigr).
 \label{eq:pre-young}
\end{align}
Young's inequality, with coefficient $g_H(1-s)/4$, gives
\[
 2D_H\bar r\sqrt{E_2}
 \leq\frac{g_H(1-s)}4\bar r^2
 +\frac{4D_H^2}{g_H(1-s)}E_2.
\]
After substitution in \eqref{eq:pre-young}, a nonnegative internal term
remains, while the coefficient of $sE_2$ is at least
\[
 a_H(1-s)^{n-1}-4D_Hs
 -\frac{4D_H^2}{g_H(1-s)}s^2.
\]
For $s\leq1/2$, Bernoulli's inequality and
$1/(1-s)\leq2$ make this at least
\[
 a_H-\left(a_H(n-1)+4D_H+\frac{8D_H^2}{g_H}\right)s.
\]
The fourth restriction in \eqref{eq:complete-core-radius} makes the last
quantity nonnegative.  This proves the theorem.  With all restrictions
strict, either $E_2>0$ or $\bar r>0$ whenever $V\ne0$, and the corresponding
retained term is strict.
\end{proof}


\section{Exact constants for all current open rows}

The graph6 strings below unambiguously specify the Atlas representatives.
Deletion--contraction computes \eqref{eq:target-remainder}.  The value $K_H$
is then fixed by \eqref{eq:constants}; the table lists $2m-n,C_H$, the
adaptive radius $\rho_H$, and the fixed cut-only radius $\sigma_H$.  The
chromatic number is also shown, so the required density threshold is explicit.

\begin{center}
\scriptsize
\setlength{\tabcolsep}{5pt}
\begin{longtable}{r@{\;}lrrrcc@{\quad}r@{\;}lrrrcc}
$H$&graph6&$\chi$&$2m-n$&$C_H$&$\rho_H$&$\sigma_H$&
$H$&graph6&$\chi$&$2m-n$&$C_H$&$\rho_H$&$\sigma_H$\\ \hline
43&\verb|Dlc|&3&7&35&$7/172$&$1/22$&118&\verb|Eht?|&3&8&76&$2/87$&$1/41$\\
122&\verb|Eld?|&3&8&76&$2/87$&$1/41$&124&\verb|Exd?|&3&8&76&$2/87$&$1/41$\\
126&\verb|ERUO|&3&8&76&$2/87$&$1/41$&127&\verb|EZEG|&3&8&82&$2/93$&$1/44$\\
130&\verb|E{CW|&3&8&64&$2/75$&$1/35$&137&\verb|EjsG|&3&10&99&$5/221$&$5/213$\\
139&\verb|Eju?|&3&10&99&$5/221$&$5/213$&145&\verb|Exe_|&3&10&117&$5/257$&$5/249$\\
147&\verb|EhMg|&3&10&117&$5/257$&$5/249$&148&\verb|EyUG|&3&10&117&$5/257$&$5/249$\\
151&\verb|EhdW|&3&10&135&$5/293$&$1/57$&152&\verb|EhNG|&3&10&117&$5/257$&$5/249$\\
153&\verb|Ehf_|&3&10&123&$5/269$&$5/261$&157&\verb|E~`_|&4&12&134&$3/146$&$3/143$\\
160&\verb|E~@g|&4&12&134&$3/146$&$3/143$&168&\verb|E^MG|&3&12&170&$3/182$&$3/179$\\
169&\verb|Exf_|&4&12&158&$3/170$&$3/167$&171&\verb|Ehew|&3&12&176&$3/188$&$3/185$\\
174&\verb|EtTg|&3&12&194&$3/206$&$3/203$&178&\verb|En}?|&4&14&181&$7/387$&$7/383$\\
181&\verb|E^mG|&4&14&205&$7/435$&$7/431$&185&\verb|Exv_|&4&14&229&$1/69$&$7/479$\\
188&\verb|EzNG|&3&14&247&$7/519$&$7/515$&194&\verb|E~wW|&4&16&300&$4/313$&$1/78$\\
196&\verb|ER~g|&4&16&300&$4/313$&$1/78$&199&\verb|El^g|&4&16&324&$4/337$&$1/84$\\
203&\verb|E~^G|&4&18&395&$9/817$&$9/817$&&&&&&&
\end{longtable}
\normalsize
\end{center}

Every adaptive radius lies between $9/817$ and $7/172$, and every fixed
cut-only radius lies between $9/817$ and $1/22$.  Since $E_2\leq1$, the
hypothesis gives $s\leq7/172$ and hence $q\leq2s<1/2$.  Therefore all covered
graphons automatically lie in the conjectured density range for the
three-chromatic rows; for the four-chromatic rows the displayed individual
radii are at most $3/146$, so $q<1/3$ there as well.

For the fully arbitrary theorem, the following table records every additional
constant in \eqref{eq:complete-core-radius}.  Here $g=g_H$,
$\Lambda=\Lambda_H$, and $N=N_3(H)$.  The radii are conservative because all
higher-rank signs and the positive $Y_V$ term are discarded.

\begin{center}
\scriptsize
\setlength{\tabcolsep}{4pt}
\begin{longtable}{r|rrrrr|c@{\quad}r|rrrrr|c}
$H$&$D_H$&$g$&$\Lambda$&$N$&$\Omega$&$\tau_H$&
$H$&$D_H$&$g$&$\Lambda$&$N$&$\Omega$&$\tau_H$\\ \hline
43&97&8&161&1&18&$7/9825$&118&241&11&481&1&37&$22/118923$\\
122&241&10&461&1&37&$5/29668$&124&241&10&464&1&37&$5/29668$\\
126&241&10&454&1&31&$5/29668$&127&266&9&454&1&31&$9/71998$\\
130&197&8&376&2&30&$8/39637$&137&316&14&1173&3&131&$35/204311$\\
139&316&13&1134&3&131&$13/81593$&145&385&14&1211&2&131&$1/8629$\\
147&385&13&1184&2&126&$13/120647$&148&385&13&1207&2&122&$13/120647$\\
151&454&13&1237&1&138&$65/836593$&152&385&12&1159&2&122&$3/30122$\\
153&410&12&1174&2&128&$3/34127$&157&435&17&2725&5&423&$17/128700$\\
160&435&16&2638&5&423&$8/64275$&168&573&17&2893&3&421&$17/222218$\\
169&529&16&2778&4&406&$8/94731$&171&598&16&2903&3&413&$2/30209$\\
174&667&16&2972&2&436&$8/150115$&178&598&20&6138&7&1213&$35/363759$\\
181&692&20&6482&6&1123&$35/485959$&185&786&20&6598&5&1193&$35/625831$\\
188&855&20&6775&4&1215&$7/147950$&194&1043&24&14648&7&3229&$48/1100605$\\
196&1043&24&14852&7&3091&$48/1100605$&199&1137&24&14992&6&3276&$16/435551$\\
203&1394&28&32316&9&8052&$63/1963067$&&&&&&&
\end{longtable}
\normalsize
\end{center}

Thus every measurable single-hub perturbation around a complete core is
nonnegative for an explicit fixed mass interval; the radii range from
$63/1963067$ to $7/9825$.  Since $q\leq2s-s^2$, these intervals lie strictly
inside the required density range.

For independent checking, the exact chromatic polynomials and targets are
grouped below.  These twenty rows cover the Atlas IDs in order of appearance
in the preceding table.

\scriptsize
\begin{longtable}{l@{\quad}p{6.6cm}p{6.6cm}}
Atlas IDs&$\chi_H(x)$&$\Phi_H(p)$\\ \hline
43&$x(x-1)(x-2)(x^2-3x+3)$&$p(2p-1)(3p^2-3p+1)$\\
118,122,124,126&$x(x-1)^2(x-2)(x^2-3x+3)$&$p^2(2p-1)(3p^2-3p+1)$\\
127&$x(x-1)(x-2)^2(x^2-2x+2)$&$p(2p-1)^2(2p^2-2p+1)$\\
130&$x(x-1)^3(x-2)^2$&$p^3(2p-1)^2$\\
137,139&$x(x-1)^2(x-2)^3$&$p^2(2p-1)^3$\\
145,147,148,152&$x(x-1)(x-2)^2(x^2-3x+3)$&$p(2p-1)^2(3p^2-3p+1)$\\
151&$x(x-1)(x-2)^2(x^2-3x+4)$&$p(2p-1)^2(4p^2-5p+2)$\\
153&$x(x-1)(x-2)(x^3-5x^2+9x-7)$&$p(2p-1)(7p^3-12p^2+8p-2)$\\
157,160&$x(x-1)^2(x-2)^2(x-3)$&$p^2(2p-1)^2(3p-2)$\\
168&$x(x-1)(x-2)^2(x^2-4x+5)$&$p(2p-1)^2(5p^2-6p+2)$\\
169&$x(x-1)(x-2)(x-3)(x^2-3x+3)$&$p(2p-1)(3p-2)(3p^2-3p+1)$\\
171&$x(x-1)(x-2)(x^3-6x^2+13x-11)$&$p(2p-1)(11p^3-20p^2+13p-3)$\\
174&$x(x-1)(x-2)(x^3-6x^2+14x-13)$&$p(2p-1)(13p^3-25p^2+17p-4)$\\
178&$x(x-1)^2(x-2)(x-3)^2$&$p^2(2p-1)(3p-2)^2$\\
181&$x(x-1)(x-2)^3(x-3)$&$p(2p-1)^3(3p-2)$\\
185&$x(x-1)(x-2)(x-3)(x^2-4x+5)$&$p(2p-1)(3p-2)(5p^2-6p+2)$\\
188&$x(x-1)(x-2)(x^3-7x^2+18x-17)$&$p(2p-1)(17p^3-33p^2+22p-5)$\\
194,196&$x(x-1)(x-2)(x-3)(x^2-5x+7)$&$p(2p-1)(3p-2)(7p^2-9p+3)$\\
199&$x(x-1)(x-2)(x-3)(x^2-5x+8)$&$p(2p-1)(3p-2)(8p^2-11p+4)$\\
203&$x(x-1)(x-2)(x-3)(x^2-6x+10)$&$p(2p-1)(3p-2)(10p^2-14p+5)$
\end{longtable}
\normalsize


\section{Endpoint, assumptions, and limitations}

At $p=1$, $V=0$ almost everywhere, hence $W=1$, $t(H,W)=1$, and monicity of
$\chi_H$ gives $\Phi_H(1)=1$.  This case is independent of the normalization
of $A$ and $B$ and uses no division by $1-p$.

The proof uses only the union bound and its exact nonnegative remainder,
Bernoulli's inequality, the scalar convexity lemma, Cauchy--Schwarz, Young's
inequality, the previously proved maximum-degree complement theorem, and the
exact chromatic polynomial.  All graphon integrations are bounded, so Tonelli
applies.  There is no regularity, step representation, approximation by
finite graphs, pointwise lower degree hypothesis, or product-of-fibre-means
assumption.  In particular, the cross fibres may cross arbitrarily and the
kernel on $A^2$ is completely unrestricted.

The exact boundary is now the complete-core hypothesis
\eqref{eq:complete-core}: the graphon must equal one on one set of measure
$1-s$.  Theorem~\ref{thm:complete-core} covers every possible kernel on all
remaining cells, uniformly for $s\leq\tau_H$.  In particular it already
covers any number of interacting hubs whose union is contained in $A$.
The independent culmination
\path{notes/arbitrary_graphon_high_density_stability.tex} removes even this
boundary: it chooses $A$ to be the high-complement-degree set of an arbitrary
graphon, applies the present quantitative margin to all defects incident to
$A$, and applies maximum-degree stability to the residual on $B^2$.  Thus the
combined theorem closes a uniform high-density interval for every graphon.
Intermediate densities remain unresolved, so no catalogue status changes.

Run
\begin{verbatim}
python codes/principled_negative_tests.py
\end{verbatim}
from the repository root.  The line headed ``arbitrary-fibre persistent-hub
stability theorem'' reconstructs $C_H,K_H,\rho_H$ for all $29$ rows from the
exact chromatic polynomial.  The crossing-cut line reconstructs every
$\sigma_H$.  The line headed ``complete-core fully arbitrary single-hub
theorem'' additionally reconstructs $D_H,g_H,\Lambda_H,N_3(H),\Omega_H$ and
all $29$ values of $\tau_H$.

For formal verification, first formalize Lemma~\ref{lem:scalar}, the exact
union-bound remainder, Lemma~\ref{lem:superadditive}, and the gain from maps
using exactly one exceptional vertex.  Reuse the maximum-degree theorem for
the internal component, then formalize the target-difference and Young bounds.
The Atlas layer needs the exact polynomial coefficients and displayed finite
constants.  Since the theorem is local in the exceptional mass, no Atlas
status module should change.

\end{document}
